%% Custom definitions for the technical content

\newcommand{\todo}[1]{\textcolor{red}{TODO: #1}\xspace}
\newcommand{\note}[1]{\textcolor{blue}{NOTE: #1}\xspace}
\newcommand{\revise}[1]{\textcolor{blue}{#1}\xspace}

\newcommand{\oursystem}{\textsc{SystemName}\xspace}

\newcommand{\paratitle}[1]{\vspace{5pt}\noindent\textbf{#1}.}
\newcommand{\paratitlenodot}[1]{\vspace{5pt}\noindent\textbf{#1}}

\newcommand*\circleb[1]{\tikz[baseline=(char.base)]{
    \node[shape=circle,fill,inner sep=.5pt](char){\textcolor{white}{\small #1}};}}
\newcommand*\circlew[1]{\tikz[baseline=(char.base)]{
    \node[shape=circle,draw,inner sep=.5pt](char){\textcolor{black}{\small #1}};}}

% --- Plain, centered captions (label handled by \caption) ---
\newcommand{\figcap}[2]{\caption{#1. #2}}  % use inside figure (below the graphic)
\newcommand{\tabcap}[2]{\caption{#1. #2}}  % use inside table (above the tabular)

% Back-compat: keep \floatcap working for FIGURES so old figure code compiles
\let\floatcap\figcap


\newcommand{\thesisfootnote}[1]{\footnote{\small #1}}

\newmdenv[  % reduce quote margins
    leftmargin=0.6em,
    rightmargin=0.6em,
    innerleftmargin=0.4em,
    innerrightmargin=0.4em,
    skipabove=1em,
    skipbelow=0.2em,
    linewidth=0pt,
]{tightquote}
